% =============================================================================
% LLM-BASED WEIGHT STRESS TEST SUBSECTION
% File: delphi_mrc_validation.tex
% NOTE: This section presents an LLM-based heuristic exploration of MRC
%       weight assignments. It is NOT a Delphi study, does NOT constitute
%       expert elicitation, and provides NO empirical validation.
%       Human-panel validation via Delphi+AHP is identified as future work.
% =============================================================================

\subsection{LLM-Based Weight Stress Test}
\label{sec:mrc_delphi}

To complement the author-assigned weights and the weight-sensitivity analysis
in Section~\ref{sec:mrc_weights}, we report an LLM-based structured
exploration of alternative weight assignments conditioned on seven
illustrative professional profiles. This is \emph{not} a Delphi study and
does not constitute expert elicitation, empirical validation, or statistical
inference from human respondents. No claim of expert consensus, scientific
validity, or validated weights is made. Human-panel validation via
Delphi+AHP with practitioner respondents is identified as future work
(Section~\ref{sec:challenges}, Gap~G7).

\begin{mdframed}[backgroundcolor=orange!6,linecolor=orange!50,linewidth=0.8pt,
  innertopmargin=5pt,innerbottommargin=5pt]
\small\textbf{Scope and Limitations.}
All responses reported below were generated by prompting a large language
model (LLM) conditioned on the stated professional profile descriptions.
The process explores how different domain perspectives \emph{might} weight
the MRC dimensions, not how they \emph{do}. LLM-generated responses do not
constitute expert opinions, do not sample from a practitioner population,
and do not support statistical inference. Any agreement statistics
reported below are computed over LLM outputs and carry no inferential
validity. This section serves as a structured heuristic stress test,
pending replacement by human-panel data.
\end{mdframed}

\subsubsection*{Illustrative Profiles}

Table~\ref{tab:delphi_personas} defines seven illustrative professional
profiles used to condition LLM responses. The profiles were designed
to span domain-specific viewpoints and stress-test whether the author
weight ordering is sensitive to professional perspective.

\begin{table*}[!t]
\caption{Illustrative Profiles Used in LLM Weight Stress Test (Not Expert Panel)}
\label{tab:delphi_personas}
\centering
\renewcommand{\arraystretch}{1.15}
\footnotesize
\begin{tabular}{@{}clp{4.5cm}@{}}
\toprule
\textbf{ID} & \textbf{Profile description} & \textbf{Primary domain lens} \\
\midrule
P1 & CISO, Fortune-500 financial institution & PKI / TLS; data-sensitivity \\
P2 & NIST PQC standardisation contributor   & Algorithm feasibility \\
P3 & ICS/OT security architect, power utility & Operational feasibility \\
P4 & IoT security researcher, embedded systems & Hardware-constrained migration \\
P5 & Cloud security engineer, hyperscaler   & TLS / VPN at scale \\
P6 & Blockchain / DLT security specialist   & Consensus-layer migration \\
P7 & Academic PQC researcher (migration focus) & Balanced/theoretical \\
\bottomrule
\end{tabular}
\end{table*}

\subsubsection*{Round~1: Simulated AHP Responses}

Each profile conditioned an LLM pairwise comparison among the three MRC
weight dimensions: Exposure ($E$), Feasibility ($F$), and Regulatory
Deadline ($D$). Table~\ref{tab:delphi_r1_ahp} reports the LLM-generated
ratios and priority vectors derived by normalised column averaging.

\begin{table*}[!t]
\caption{Round~1 Simulated AHP Comparisons and Priority Vectors (LLM-generated)}
\label{tab:delphi_r1_ahp}
\centering
\renewcommand{\arraystretch}{1.2}
\footnotesize
\begin{tabular}{@{}ccccccccc@{}}
\toprule
\textbf{ID} & $a_{EF}$ & $a_{ED}$ & $a_{FD}$ & $w_E$ & $w_F$ & $w_D$ & CR \\
\midrule
P1 & 4     & 2    & 1/3  & 0.557 & 0.123 & 0.320 & 0.016 \\
P2 & 2     & 4    & 3    & 0.557 & 0.320 & 0.123 & 0.016 \\
P3 & 1/2   & 3    & 5    & 0.309 & 0.581 & 0.110 & 0.003 \\
P4 & 1/4   & 2    & 5    & 0.201 & 0.681 & 0.118 & 0.021 \\
P5 & 5     & 3    & 1/2  & 0.648 & 0.122 & 0.230 & 0.003 \\
P6 & 1/3   & 3    & 5    & 0.261 & 0.633 & 0.106 & 0.033 \\
P7 & 2     & 3    & 2    & 0.539 & 0.297 & 0.164 & 0.008 \\
\midrule
\multicolumn{4}{@{}l}{\textbf{Aggregate (geometric mean, LLM)}} &
  \textbf{0.459} & \textbf{0.367} & \textbf{0.174} & \textbf{0.003} \\
\bottomrule
\multicolumn{8}{@{}l}{\scriptsize $a_{XY}$: importance of $X$ relative to $Y$ (AHP 9-pt scale).
  CR: consistency ratio. All values LLM-generated; not empirically collected.}
\end{tabular}
\end{table*}

\noindent\textbf{Observations.}
\begin{itemize}
  \item \textbf{Exposure weight broadly consistent with author value.} The
  aggregate yields $w_E{=}0.459$, within 2\% of the author value (0.45).
  Domain-specialists constrained by hardware (P3, P4, P6) produce higher
  feasibility weights, reflecting the domain insight that urgency becomes
  moot when migration is technically infeasible.

  \item \textbf{Feasibility weight slightly elevated.} The aggregate
  $w_F{=}0.367$ exceeds the author's 0.35, driven by the three
  constrained-environment profiles (P3, P4, P6). This suggests the author
  weight may underweight operational constraints in IoT, ICS/OT, and
  blockchain domains.

  \item \textbf{Regulatory deadline weight lower.} The aggregate
  $w_D{=}0.174$ falls below the author's 0.20, consistent across
  most profiles. This aligns with the sensitivity analysis (Section~%
  \ref{sec:mrc_weights}): deadline scores have limited tier-determining
  power given their narrow range in the reference profiles.
\end{itemize}

\subsubsection*{Round~2: Simulated Revision}

A second LLM pass was conditioned on Round~1 aggregate weights to check
directional stability. Table~\ref{tab:delphi_r2_ahp} reports the revised
values.

\begin{table*}[!t]
\caption{Round~2 Simulated AHP Comparisons and Priority Vectors (LLM-generated)}
\label{tab:delphi_r2_ahp}
\centering
\renewcommand{\arraystretch}{1.2}
\footnotesize
\begin{tabular}{@{}ccccccccl@{}}
\toprule
\textbf{ID} & $a_{EF}$ & $a_{ED}$ & $a_{FD}$ & $w_E$ & $w_F$ & $w_D$ & CR &
  \textbf{Change vs R1} \\
\midrule
P1 & 3     & 2    & 1/3  & 0.525 & 0.142 & 0.334 & 0.046 & $a_{EF}$: 4$\to$3 \\
P2 & 2     & 4    & 3    & 0.557 & 0.320 & 0.123 & 0.016 & Unchanged \\
P3 & 2/3   & 3    & 4    & 0.359 & 0.517 & 0.124 & 0.001 & $a_{EF}$: $\tfrac{1}{2}{\to}\tfrac{2}{3}$; $a_{FD}$: 5$\to$4 \\
P4 & 1/3   & 2    & 4    & 0.240 & 0.623 & 0.137 & 0.016 & $a_{EF}$: $\tfrac{1}{4}{\to}\tfrac{1}{3}$; $a_{FD}$: 5$\to$4 \\
P5 & 4     & 3    & 1/2  & 0.623 & 0.137 & 0.240 & 0.016 & $a_{EF}$: 5$\to$4 \\
P6 & 1/2   & 3    & 4    & 0.320 & 0.557 & 0.123 & 0.016 & $a_{EF}$: $\tfrac{1}{3}{\to}\tfrac{1}{2}$; $a_{FD}$: 5$\to$4 \\
P7 & 2     & 3    & 2    & 0.539 & 0.297 & 0.164 & 0.008 & Unchanged \\
\midrule
\multicolumn{4}{@{}l}{\textbf{Aggregate (geometric mean, LLM)}} &
  \textbf{0.471} & \textbf{0.348} & \textbf{0.181} & \textbf{0.004} \\
\bottomrule
\end{tabular}
\end{table*}

The aggregate ordering $w_E{>}w_F{>}w_D$ is stable across both rounds.
The Round~2 aggregate $(0.471, 0.348, 0.181)$ lies within 2.1 percentage
points of the author values on each dimension.

\subsubsection*{LLM-Suggested Sub-score Adjustments}

The LLM outputs also flagged three domains where the author sub-scores
may warrant revision based on domain-specific reasoning:

\begin{enumerate}
  \item \textbf{IoT/Embedded $F$: possible overestimate.} Profiles P4 and
  P2 suggested $F{=}0.500$ overstates feasibility for the sub-Cortex-M4
  deployed base, proposing $F{\approx}0.400$ based on ROM/stack constraints
  for CRYSTALS-Kyber-512 on Cortex-M0/M0+.

  \item \textbf{ICS/OT $F$: possible overestimate.} Profile P3 suggested
  $F{=}0.500$ underweights 15--25 year equipment lifecycles and IEC~62443
  change-control requirements; $F{\approx}0.430$ was proposed.

  \item \textbf{Blockchain $F$: possible overestimate.} Profile P6 argued
  that consensus-layer hard-fork coordination reduces effective feasibility;
  $F{\approx}0.300$ was proposed, with $E$ slightly raised to 0.630 to
  reflect on-chain ECDSA public key exposure.
\end{enumerate}

These LLM-generated suggestions serve as prompts for future human-panel
validation rather than calibrated corrections. The full sensitivity
analysis in Table~\ref{tab:mrc_sensitivity} covers the range of sub-score
variation that would be required to change any tier assignment.

\subsubsection*{Tier Stability Under Alternative Weights}

Table~\ref{tab:delphi_tier_stability} compares tier assignments under
four scenarios combining author vs.\ LLM-aggregate weights with original
vs.\ LLM-suggested sub-scores.

\begin{table*}[!htb]
\caption{Tier Stability Under Alternative Weights and Sub-scores (LLM-generated scenarios).
  MRC$=$0.45$E$+0.35(1$-F$)+0.20(1$-D$) (author) or
  0.471$E$+0.348(1$-F$)+0.181(1$-D$) (LLM aggregate).
  Thresholds: T1$\geq$0.75, T2$\geq$0.50, T3$\geq$0.25, T4$<$0.25.}
\label{tab:delphi_tier_stability}
\centering
\renewcommand{\arraystretch}{1.2}
\footnotesize
\begin{tabular}{@{}lccccl@{}}
\toprule
 & \multicolumn{2}{c}{\textbf{Original sub-scores}}
 & \multicolumn{2}{c}{\textbf{LLM-adjusted sub-scores}} \\
\cmidrule(lr){2-3}\cmidrule(lr){4-5}
\textbf{Domain}
  & \shortstack{Author\\weights} & \shortstack{LLM\\weights}
  & \shortstack{Author\\weights} & \shortstack{LLM\\weights}
  & \textbf{Verdict} \\
\midrule
TLS/Web      & T2 (0.703) & T2 (0.696) & T2 (0.691) & T2 (0.683) & Stable T2 \\
PKI/Certs    & T1 (0.788) & T1 (0.787) & T1 (0.799) & T1 (0.798) & Stable T1 \\
IoT/Embedded & T2 (0.605) & T2 (0.601) & T2 (0.640) & T2 (0.636) & Stable T2 \\
ICS/OT       & T2 (0.613) & T2 (0.618) & T2 (0.627) & T2 (0.633) & Stable T2 \\
Blockchain   & T2 (0.640) & T2 (0.636) & T2 (0.689) & T2 (0.685) & Stable T2 \\
VPN/ZTNA     & T1 (0.888) & T1 (0.882) & T1 (0.881) & T1 (0.875) & Stable T1 \\
\bottomrule
\multicolumn{6}{@{}l}{\scriptsize MRC scores in parentheses. LLM-generated scenarios; not
  empirically validated.}
\end{tabular}
\end{table*}

\subsubsection*{Summary}

\begin{mdframed}[backgroundcolor=blue!4,linecolor=blue!40,linewidth=0.6pt,
  innertopmargin=4pt,innerbottommargin=4pt]
\small\textbf{LLM Weight Stress Test: Key Findings.}
\begin{enumerate}[leftmargin=*,itemsep=2pt,topsep=2pt]
  \item \textbf{Weight ordering robust across profiles.}
  The aggregate ordering $w_E{>}w_F{>}w_D$ is stable across both
  LLM simulation rounds. Aggregate weights $(0.471, 0.348, 0.181)$ lie
  within 2.1\,pp of author values.

  \item \textbf{Domain-specific feasibility may be overstated.}
  IoT, ICS/OT, and Blockchain feasibility scores are flagged for
  potential downward revision; these require primary-source validation
  before recalibration.

  \item \textbf{Tier assignments fully stable.}
  All six domains retain their tier across all four scenarios.
  PKI and VPN/ZTNA remain T1; all others remain T2.

  \item \textbf{Human validation required.}
  These findings are LLM-generated heuristic prompts, not validated
  weights. A practitioner Delphi survey is identified as Gap~G7
  (Section~\ref{sec:challenges}).
\end{enumerate}
\end{mdframed}
